% !TEX program = pdflatex
% =====================================================================
% ICCSCI 2026 — Procedia Computer Science
% SLR Paper: Machine Learning for Village Fund Anomaly Detection:
%   A Systematic Literature Review of the Operationalization Chasm
%   in Indonesian Anti-Corruption Information Systems
%
% Template: elsarticle + ecrc (Elsevier Procedia CRC)
% Citation style: Vancouver numbered [#]
% =====================================================================

\documentclass[3p,times,procedia]{elsarticle}
\flushbottom

%% ecrc package — required for Procedia CRC layout
\PassOptionsToPackage{procedia}{ecrc}
\usepackage{ecrc}
\usepackage[bookmarks=false]{hyperref}
    \hypersetup{colorlinks,
      linkcolor=blue,
      citecolor=blue,
      urlcolor=blue}

%% Core packages
\usepackage{graphicx}
\usepackage{stfloats}
\usepackage{amsmath}
\usepackage{amssymb}
\usepackage{booktabs}
\usepackage{array}
\usepackage{tabularx}
\usepackage{etoolbox}

\graphicspath{{../../analysis/bibliometric/}{../../analysis/themes/charts/}}

%% Journal metadata
\volume{00}
\firstpage{1}
\journalname{Procedia Computer Science}
\runauth{[BLINDED FOR REVIEW]}
\jid{procs}

% =====================================================================
\begin{document}
\begin{frontmatter}

\dochead{The 13th International Conference on Computer Science and Computational Intelligence (ICCSCI 2026)}

\title{A Systematic Review of the Literature on The Operationalisation Chasm in Indonesian Anti-Corruption Information Systems for Machine Learning to Find Village Fund Anomalies}

\author[a]{Yodihartomo, Farrell\corref{cor1}}
\author[b]{Faiq, Humam}
\address[a]{Bina Nusantara University}
\address[b]{Komisi Pemberantasan Korupsi}
\cortext[cor1]{Corresponding author.}

\begin{abstract}
The issue of village fund corruption in Indonesia is an ongoing governance failure. Since 2015 until 2023, KPK identified more than 600 Dana Desa violations with total financial loss exceeding IDR 433 billion. While advances in machine learning-based fraud detection continue, none of them have so far been applied to the particular context of Indonesian village fund governance. Thus, in order to develop a solution that could leverage the existing advancements in the field, this systematic literature review seeks to examine: (i) which machine learning approaches have been developed, (ii) which feature-engineering strategies have been tested, and (iii) what obstacles there are to their implementation within the scope of Dana Desa anomaly detection. The PRISMA 2020 protocol was used for conducting this literature review with 1,001 sources pre-filtered and 45 relevant papers from 2018 to 2026 included in the final sample. Thematic analysis, bibliometrics, and matrix mapping based on Design Science Research (DSR) framework allowed to identify four major research themes as follows: (a) Operationalization Chasm – lack of ML-detection-village-governance connection; (b) Scalability Illusion – performance metrics assume inappropriate circumstances; (c) Absence of Information Systems theory in 62\% of technical papers; and (d) Ground Truth Paradox – circular validation via synthetic datasets. Three research gaps inform an original empirical DSR artefact.
\end{abstract}

\begin{keyword}
village fund corruption \sep machine learning \sep anomaly detection \sep
systematic literature review \sep information systems governance
\end{keyword}

\end{frontmatter}
\email{[email-blinded-for-review@institution.edu]}

\vspace*{-6pt}

% =====================================================================
\section{Introduction}
\label{sec:intro}
% =====================================================================

The village funds administration in Indonesia exists at such a significant scale that it makes the administration simultaneously vital and susceptible to abuse. Through the Villages Fund law (UU No. 6/2014), the country allocates more than IDR 70 trillion per fiscal year to more than 74,000 village governments. Fiscal decentralization model involves allocating money directly into the budget of the village governments without the oversight of independent auditors; the process allows for the involvement of village officials into the corrupt practices. Based on the Corruption Eradication Commission (KPK)'s findings between 2015 and 2023 \cite{method_kpk2023}, the village funds corruption includes misappropriation of money, fictitious project spending, and fraudulent budget absorption. The financial vulnerability is systematic because the pattern of recurrence persists during several consecutive years, demonstrating the existence of structural principal-agent problems \cite{method_jensen1976}.

At present, most of the governance measures applied by the government focus on ex post audit cycles identifying fraud after the disbursements of the budget have occurred. The capacity constraints involved are structural because only a small portion of the villages undergo audits performed by the Inspektorat Daerah (APIP) every year. In addition, the audit conducted by the Badan Pemeriksa Keuangan (BPK) covers only aggregate financial reports, without delving deeper into the expenditure data. As a result, there is no systematic real-time screening of the data stored in the information system Siskeudes. The literature on ML-based financial fraud detection has grown extensively; however, the studies have focused on commercial databases and have reported high AUC-ROC performance of Isolation Forest, Local Outlier Factor (LOF), GNNs, and autoencoder algorithms \cite{p003,p007,p015,p022,p026,p041,p043}. Nevertheless, it remains unclear whether such solutions could be applied to village governance under fragmented records, unlabeled data, and subnational differences.

These three research questions inform the research objectives of this Systematic Literature Review (SLR):
\begin{itemize}
\item \textbf{RQ1}: What ML-based anomaly detection approaches are applicable in the field of finance, and how do they perform under different conditions?
\item \textbf{RQ2}: What are the feature engineering methods for modeling corruption as computational indicators, and to what extent can they be used in a public-sector context?
\item \textbf{RQ3}: How are ML-based approaches restricted in their applicability to decentralized village fund governance, and what design specifications need to be considered for primary research?
\end{itemize}

From an information systems (IS) point of view, the SLR defines ML anomaly detection as a means of IS governance, specifically an artefact that needs to be implemented in institutions in order to serve as a tool for human decision-making \cite{method_hevner2004}. A holistic evaluation of such an artefact takes into account factors such as predictive accuracy, institutional applicability, interpretability, and coherence with theories of the principal–agent relationship as an explanation for corruption \cite{method_jensen1976}. The contribution of the SLR is to provide a structured body of knowledge regarding the design of such an artefact, as well as validation requirements.


% =====================================================================
\section{Related Work}
\label{sec:related}
% =====================================================================

\subsection{ML-Based Financial Fraud Detection}

According to contemporary literature reviews, there are three methodological families: ensembles of decision trees, deep learning neural network models (including LSTM, Autoencoder, and Transformer), and graph models \cite{p001,p003,p004,p022,p037}. Unsupervised learning algorithms have become popular due to the lack of ground truth labels: Isolation Forest \cite{p015}, LOF \cite{p041}, and One-Class SVM \cite{p015,p041}. The point is that the presence or absence of labels determines the applicability of the respective method family. Thus, in the Dana Desa setting where fraud convictions exist only as isolated cases in KPK investigations and not as labeled dataset, unsupervised or semi-supervised preprocessing would be necessary for using existing supervised machine learning models for high AUC-ROC values \cite{p003,p007,p022,p026,p039}.

Another issue with existing systematic literature reviews (SLRs) concerns the fact that performance evaluation occurs for the private sector settings where data comes from credit cards transactions, insurance claims, and AML networks, but then conclusions regarding applicability to public administration scenarios are drawn despite the lack of validation concerning the same institutional conditions \cite{p001,p022,p037}. It is an example of the Scalability Illusion when the model works well under one institutional context and does not work because of that context in another. Even procurement fraud detection \cite{p008,p043}, which can be considered most similar to the village fund expenditure anomaly, is applied to centralized government procurement processes.

\subsection{Village Fund Governance and Corruption Patterns}

Literature on governance through Dana Desa provides a well-established corruption typology grounded in agency theory \cite{method_jensen1976}, fraud triangles, and institutional control mechanisms. Some of the established practices include budget absorption frauds, fictitious procurement processes, and administrative frauds \cite{method_kpk2023}. In terms of the governance approach, institutional settings of weak monitoring, rotating village chiefs, and insufficient audit capabilities provide the necessary context to transform latent incentives into observable instances of corruption. While the digital governance field \cite{p013,p016,p039,p044} has contributed with information system perspectives on the technological dimensions of the corruption phenomena under discussion, there is no computational method available to detect these forms of corruption. The problem demonstrates the separation of scholarly communities on the one hand, since the governance research community knows most about fraud mechanisms, while on the other hand the machine learning community can create detection tools but does not utilize them because of different fraud mechanics in governance vs business cases.

\subsection{Theoretical Framework: DSR and IS Success}

In the current study, the three-cycle DSR methodology of Hevner et al. (2004) \cite{method_hevner2004} serves as the primary tool of analysis. According to DSR criteria, an artifact that fulfills Rigor and Relevance but does not pass the DESIGN cycle is unable to contribute any knowledge to the IS domain. As can be seen from the DSR framework matrix depicted in Section \ref{sec:results}, no one paper passes the DESIGN × village criteria.

Moreover, the DeLone and McLean IS success model (2003) \cite{method_delone2003} complements DSR assessment through providing other artefact-assessment criteria besides AUC-ROC: quality of information (the anomaly score should be interpretable), quality of system (integration with the Siskeudes workflow is required), and net benefits (should alleviate the APIP workload). Finally, the Agency Theory by Jensen \& Meckling (1976) \cite{method_jensen1976} explains the principal-agent corruption model, where the structure variable targeted for elimination is the information asymmetry of the principal.

% =====================================================================
\section{Methodology}
\label{sec:method}
% =====================================================================

This systematic review was conducted based on the PRISMA 2020 guideline \cite{method_prisma2020} and using a thematic synthesis approach introduced by Thomas and Harden \cite{method_thomas2008}. Additional evidence came from bibliometrics \cite{method_donthu2021} and a design science research (DSR) framework matrix \cite{method_hevner2004}. Transparency in the reporting of SLR follows Templier and Paré (2018) \cite{method_templier2018}. The screening process flow in PRISMA 2020 is illustrated in Figure~\ref{fig:framework}.

\begin{figure}[htb]
\centering
\includegraphics[width=0.8\columnwidth]{./charts/Fig6_prisma_flow.png}
\caption{PRISMA 2020 screening flow: database identification, title/abstract screening, quality assessment, domain-override adjudication, and final inclusion (N\,=\,45).}
\label{fig:framework}
\end{figure}

\subsection{Search Strategy and Eligibility}

Searches were carried out in five databases: OpenAlex (via API automation), Scopus, IEEE Xplore, Web of Science, and Semantic Scholar. Two sets of keywords were used for searches: the first focused on machine learning (ML) techniques related to public financial anomalies or fraud (for RQ1 and RQ2) and the second one covered governance of village funds (for RQ3). The searches initially returned 1,001 records.

The inclusion criteria were as follows: (IC-01) the use of ML/AI algorithms for detecting financial anomalies or frauds; (IC-02) any research related to public sector governance or corruption; (IC-03) IS theoretical constructs employed; (IC-04) published from 2010 to 2026; (IC-05) articles published via peer-review process; and (IC-06) written in English or Indonesian language. A \textit{domain-override protocol} retained 12 borderline papers (quality score 4.0-4.99) with direct Dana Desa relevance.

\subsection{Quality Assessment and Selection}

A composite quality rating from 0 to 10 was determined for each article based on five categories: methodological rigour, theoretical basis, contextual relevance, replicability, and level of evidence. The enrichment of the journal SCImago quartile was considered for 194 articles. The minimum criteria for inclusion were $\geq 5.0$, while the domain override criteria were $\geq 4.0$ with validated Dana Desa relevance.

Stage 1 inter-coder reliability reached $\kappa = 0.307$ initially; after adjudication, $\kappa = 1.000$. Final included corpus: \textbf{N\,=\,45} \cite{p001,p002,p003,p004,p005,p006,p007,p008,p009,p010,p011,p012,p013,p014,p015,p016,p017,p018,p019,p020,p021,p022,p023,p024,p025,p026,p027,p028,p029,p030,p031,p032,p033,p034,p035,p036,p037,p038,p039,p040,p041,p042,p043,p044,p045}.

\subsection{Synthesis Methods}

Three analytical tools were used simultaneously. The first analytical tool used was thematic synthesis \cite{method_thomas2008}. It involved the analysis of 613 coded items from 45 studies through open coding to provide 10 descriptive themes (DT1-DT10) and four analytical themes (AT1-AT4). 121 relations between the papers were identified as being either converging, extending, contradicting, silencing, or bridging. DSR research framework matrix \cite{method_hevner2004} was used to map all papers into three research cycles (DESIGN, RELEVANCE, RIGOR) across six governance contexts; empty squares indicate structure gaps.Bibliometric analysis \cite{method_donthu2021} examined temporal trends, journal distribution, and keyword co-occurrence clustering.

\subsection{Sensitivity Analysis}

Three quality thresholds were considered: T1 ($\geq$4.0, N = 45), T2 ($\geq$4.5, N = 23), and T3 ($\geq$5.0, N = 14) \cite{method_higgins2011}. T1 acts as the default threshold for reporting purposes, covering all cases where there are any overrides within the domains. Table \ref{tab:sensitivity} presents the robustness of the findings. Four themes of analysis (AT1-AT4) and the DESIGN × village interaction effect empty cell are constant, suggesting that the Chasm is not an artifact of selection criteria.

\begin{table}[htb]
\caption{Core Finding Stability Across Quality Thresholds}
\label{tab:sensitivity}
\small
\begin{tabular*}{\hsize}{@{\extracolsep{\fill}}lp{1.3cm}p{1.3cm}p{1.3cm}@{}}
\toprule
\textbf{Core Finding} & \textbf{T1 (N=45)} & \textbf{T2 (N=23)} & \textbf{T3 (N=14)} \\
\colrule
AT1: Operationalization Chasm & Stable & Stable & Stable \\
AT2: Scalability Illusion & Stable & Stable & Stable \\
AT3: IS Theory Absence & Stable & Stable & Stable \\
AT4: Ground Truth Paradox & Stable & Stable & Stable \\
DESIGN$\times$village\,=\,0 & Confirmed & Confirmed & Confirmed \\
DT5 Village governance & Stable & Unstable & Unstable \\
\botrule
\end{tabular*}
\end{table}

% =====================================================================
\section{Results}
\label{sec:results}
% =====================================================================

\subsection{Bibliometric Characteristics}

The dataset consists of 45 articles from 2018 to 2026, including one article in 2021, six in 2022, nine in 2023, twelve in 2024, fifteen in 2025, and two in 2026, accounting for 98\% of the articles published from 2022 to 2026. Articles are mostly published in IEEE Access (Q1) and Applied Sciences (Q2) under the machine learning detection tradition, while publications on the governance tradition appear mainly in regional information systems and public administration journals, some of which have no ranking in Scopus. Figure \ref{fig:trend} illustrates the distribution of the publication time frame, and Figure \ref{fig:cluster} shows the segregation of the bibliometric clusters.

\begin{figure*}[htb]
\centering
\begin{minipage}[t]{0.47\textwidth}
\centering
\includegraphics[width=\textwidth]{./charts/publication_trend.png}
\caption{Annual publication count 2021--2026 with cumulative overlay (N\,=\,45).}
\label{fig:trend}
\end{minipage}
\hfill
\begin{minipage}[t]{0.47\textwidth}
\centering
\includegraphics[width=\textwidth]{./charts/cluster_separation.png}
\caption{Bibliometric cluster separation: ML\_DETECTION vs IS\_GOVERNANCE. Zero cross-cluster keyword co-occurrence $>$2 constitutes direct evidence of the Operationalization Chasm.}
\label{fig:cluster}
\end{minipage}
\end{figure*}

The application of bibliometric clustering using a keyword fraction threshold higher than 0.6 results in a structurally segregated set of documents containing two clusters (refer to Table \ref{tab:clusters}). Also, the fact that there is no occurrence of an ML cluster keyword co-occurring more than twice with a governance cluster keyword implies that there is bibliometric proof of existence of Operationalization Chasm. The three bridging papers \cite{p008,p013,p043} serve as the only citation bridge between the two groups; however, the existence of Operationalization Chasm cannot be considered absolute due to just three papers serving as bridges.

\begin{table*}[htb]
\caption{(a) Bibliometric cluster structure (N\,=\,45).\quad (b) DSR cycle $\times$ context level matrix.}
\label{tab:clusters}
\small
\begin{minipage}[t]{0.42\textwidth}
\begin{tabular*}{\hsize}{@{\extracolsep{\fill}}lrp{3.2cm}@{}}
\toprule
\textbf{Cluster} & \textbf{N} & \textbf{Characterization} \\
\colrule
ML\_DETECTION  & 21 & Deep learning, IF, LOF; private/synthetic data; no IS theory \\
IS\_GOVERNANCE & 21 & Dana Desa, agency theory, institutional controls \\
BRIDGING       & 3  & Mixed signals --- partial overlap of both clusters \\
\botrule
\end{tabular*}
\end{minipage}
\hfill
\begin{minipage}[t]{0.53\textwidth}
\begin{tabular*}{\hsize}{@{\extracolsep{\fill}}lrrrrr@{}}
\toprule
\textbf{DSR Cycle} & \textbf{Village} & \textbf{Sub-nat.} & \textbf{Nat.} & \textbf{XN-Dev.} & \textbf{XN-Dvlp.} \\
\colrule
DESIGN    & \textbf{0} & \textbf{0} & 3  & 5  & 2 \\
RELEVANCE & 14         & 0          & 5  & 0  & 0 \\
RIGOR     & 11         & 1          & 7  & 8  & 5 \\
\colrule
\textit{Total} & 14 & 2 & 8 & 8 & 5 \\
\botrule
\end{tabular*}
\end{minipage}
\end{table*}

\subsection{Descriptive Themes}

Theme generation from 613 codes from 45 articles resulted in 10 descriptive themes (Table ~\ref{tab:dt}). The most important observation across all the research articles is the fact that 28 out of 45 articles (62\%) do not have any theory related to information systems. Looking at the breakdown based on research questions (RQ), we see an imbalance. While DT3 (22 articles) addresses RQ1 through supervised learning, all of which make the assumption of label availability that breaks down in the context of village governance, DT2 (unsupervised, 18 articles) makes up the minority.

\begin{table*}[htb]
\caption{Descriptive Theme Summary (DT1--DT10)}
\label{tab:dt}
\small
\begin{tabular*}{\hsize}{@{\extracolsep{\fill}}clrr@{}}
\toprule
\textbf{ID} & \textbf{Theme Label} & \textbf{N} & \textbf{RQ} \\
\colrule
DT1  & Corruption as computational signals    & 16 & 2   \\
DT2  & Unsupervised ML anomaly detection      & 18 & 1   \\
DT3  & Supervised detection (private sector)  & 22 & 1   \\
DT4  & Label scarcity / ground truth gap      & 19 & 3   \\
DT5  & Village fund governance \& corruption  & 26 & 2,3 \\
DT6  & IS-theoretical framing                 & 17 & 1,3 \\
DT7  & Real-time detection gap                & 4  & 3   \\
DT8  & Developing-country constraints         & 18 & 3   \\
DT9  & Graph / network-based methods          & 20 & 1   \\
DT10 & Explainability \& audit requirements   & 7  & 1,3 \\
\botrule
\end{tabular*}
\end{table*}

\begin{figure*}[htb]
\centering
\begin{minipage}[t]{0.47\textwidth}
\centering
\includegraphics[width=\textwidth]{./charts/Fig2_rq_theme_heatmap.png}
\caption{Theme $\times$ RQ coverage matrix (cell = paper count). DT7--RQ3 ($n=4$) marks the sole silencing gap for real-time deployment.}
\label{fig:heatmap}
\end{minipage}
\hfill
\begin{minipage}[t]{0.47\textwidth}
\centering
\includegraphics[width=\textwidth]{./charts/Fig4_evidence_density_rq.png}
\caption{Evidence density per RQ. Red bars = thin-evidence themes ($\leq$6 papers). DT7 in RQ3 is the sole silencing gap.}
\label{fig:evdensity}
\end{minipage}
\end{figure*}

\subsection{Analytical Themes}

Cross-paper relationship analysis produced four analytical themes (Table~\ref{tab:clusters}).

\textbf{AT1 --- The Operationalization Chasm}: There is a division of literature into two non-communicating traditions. 23 papers present ML method development without involvement of governance, and 26 papers study corruption in villages without utilizing ML techniques. There are three papers bridging the divide, but they fail to bridge two clusters. Six silencing relations confirm the division, such as the absence of ML-based studies for Dana Desa, failure to propose village fund artefacts for ML, no mapping corruption taxonomy to detectable indicators.

\textbf{AT2 --- The Scalability Illusion}: Papers from the technical cluster use unreal assumptions about data (centralized database - 4 papers, labeled data - 12 papers, and a large number of transactions). Papers from the cluster utilize synthetic validation five times, hence making circular validation possible.

\textbf{AT3 --- The Absence of IS Theory}: Only 62\% (27 out of 45) use IS theory. IS theory can be found solely in the governance cluster, and never together with detection methodology. Moreover, no papers use evaluation of any detection artefact using the DeLone–McLean IS success model or any other IS adoption framework \cite{method_delone2003}.

\textbf{AT4 --- The Ground Truth Paradox}: Five papers claim a high value of AUC-ROC on either synthetic or self-labeled datasets while admitting the lack of ground truth in the field.

\subsection{DSR Framework Matrix}

The right-hand side of Table~\ref{tab:clusters} shows which cluster each of the papers belongs to, based on their corresponding DSR cycle and governance context level. Interestingly, the \textbf{DESIGN $\times$ village cell has no papers}—an observation that stands out structurally. There has never been a paper that designs and tests the effectiveness of machine learning artefacts based on village-level government expenditure. RELEVANCE$\times$village = 14 implies that researchers have clearly defined the problem space for the study of governance. Meanwhile, RIGOR$\times$village = 11 suggests that there are frameworks available for doing rigorous research. However, what is missing is a Design-cycle artefact connecting the two.

From Figure~\ref{fig:keywords}, it is apparent how the lexical mechanism has produced the structural discrepancy. The keywords dominating ML\_DETECTION (precision, recall, AUC-ROC, ensemble, graph neural network) and IS\_GOVERNANCE (dana desa, principal-agent, inspektorat, APIP, transparency) do not co-occur above the value of 2. The reason behind the failure of the Design cycle in the village-context lies in this semantic difference.

\begin{figure}[htb]
\centering
\includegraphics[width=0.58\columnwidth]{./charts/keyword_frequency.png}
\caption{Top keyword frequencies by cluster (ML\_DETECTION vs IS\_GOVERNANCE). Zero cross-cluster co-occurrence confirms the two traditions operate in non-overlapping semantic spaces, corroborating the DESIGN$\times$village\,=\,0 cell in Table~\ref{tab:clusters} (right).}
\label{fig:keywords}
\end{figure}

\subsection{Gap Matrix}

\begin{table}[htb]
\caption{Integrated Research Gap Matrix}
\label{tab:gaps}
\small
\begin{tabular*}{\hsize}{@{\extracolsep{\fill}}clrp{7cm}@{}}
\toprule
\textbf{ID} & \textbf{Severity} & \textbf{N} & \textbf{Gap Statement} \\
\colrule
G1 & CRITICAL     & 14 & No ML artefact applies to Dana Desa financial data (DESIGN $\times$ village = 0) \\
G2 & CRITICAL     & 19 & Label scarcity eliminates supervised approaches; unsupervised methods remain methodologically underspecified \\
G3 & CRITICAL     & 40 & No validated feature engineering framework exists for village fund fraud \\
G4 & PARTIAL      & 37 & IS theory absent from 62\% of detection papers; no D\&M IS Success Model application found \\
G5 & METHODOLOGICAL & 7 & Real-time/streaming deployment architectures absent; explainability insufficient \\
\botrule
\end{tabular*}
\end{table}

The gaps G1, G2, and G3 are interrelated from a structural perspective: the gap G3 (lack of a theoretical feature set) is a precondition for the gap G1 (lack of an application) because deploying the ML construct on the data collected in villages without a theoretical foundation in a feature set would generate anomaly scores that do not have a basis for interpretation by the auditors.

% =====================================================================
\section{Discussion}
\label{sec:discussion}
% =====================================================================

\subsection{Answering the Research Questions}

\textbf{RQ1} (ML methods): Isolation Forest \cite{p015}, Local Outlier Factor (LOF) \cite{p041}, and Autoencoders \cite{p003,p015}show the best suitability for unlabeled, small-volume village fund datasets. Supervised techniques prevail in the existing corpus \cite{p003,p007,p022,p026,p039} but need labeled samples that cannot be accessed in the Dana Desa setting. Graph Neural Networks (GNNs) show better relational modeling capabilities \cite{p014,p045} than alternatives but need a transaction graph architecture that is missing in the existing village data.

\textbf{RQ2} (Feature engineering): DT1 (16 articles) reveals five transferable feature classes: single-bidder ratios, amendment rates, price increments, timing clusters, and vendor diversities \cite{p008,p043}. Two unique village features come out from DT5: budget utilization irregularities and payment timing inconsistencies \cite{method_kpk2023}. Prior research does not combine these variables into a consistent village fund detection model.

\textbf{RQ3} (Structural gaps): Three structural limitations prevent direct applicability (G1, G2, G3), corroborated by convergent triangulation of thematic analysis, bibliometric disconnection, and Design Science Research (DSR) empty spaces. The Operationalization Gap represents an epistemological gap between traditions that neither has been able to solve.

\subsection{Theoretical Contributions}

AT3 suggests that the lack of IS theory accounts for 62\% of the lack of information systems theory is not only an academic issue since the ML detection literature optimizes the AUC-ROC measure, which is unrelated to the three critical components for success (information quality, system quality, and net benefits \cite{method_delone2003}) of a detection system as a governance tool. With no grounding in IS theory, there is no design for the institutional process and the form in which audits must be explained in order to use the outputs generated by detection.

AT1 highlights an even more serious conceptual flaw in both approaches in that neither tradition provides a clear understanding of what corruption entails in terms of computation in the village fund governance process. This requires developing a corruption classification framework that is sociologically sound (in accordance with the methods employed by the Corruption Eradication Commission (KPK)).

\subsection{Implications for the Field}

First, there is a fundamental disjunction between ML techniques for detecting fraud and the requirements of governance in the public sector. The ML paradigm, which relies heavily on supervised learning from labeled data, creates an \textit{illusion of generalizability}: high AUC-ROC scores obtained under controlled conditions do not carry over into situations where ground truth cannot be institutionalized. In the case of village fund governance, the state-of-the-art methods of detection are precisely those that are least implementable.

Second, the bibliometric distance between the ML and IS governance knowledge domains goes beyond the simple lack of citations. Each body of knowledge represents distinct epistemic communities. ML specialists aim at maximizing prediction accuracy, while governance theorists focus on institutional dysfunction. In the absence of a mediating literature, each community amplifies the assumptions that are made by the other. Thus, ML scientists posit conditions that are not met in reality, while governance researchers describe problems without proposing computational solutions.

The absence of IS theories in 62\% of the reviewed papers deserves special attention. Implementing detection technologies without a theory of adoption, explanation, and institutionalization may lead to reducing political problems to technical issues. The DeLone \& McLean IS Success Model can be helpful in addressing this problem. Anomaly detection systems must be evaluated not only based on the F1 score, but also their potential to decrease information asymmetry between principals and agents.

\subsection{Limitations}

The review does not include any seminal publications that date back before 2010. The inclusion criteria for the 45 articles take into account the limitations of the database used and the search string. The quality assessment criteria discriminate against studies that examine Indonesian village administration that were published in journals with no ranking system.

% =====================================================================
\section{Conclusion}
\label{sec:conclusion}
% =====================================================================

This systematic literature review (SLR) synthesises 45 studies and identifies structural gaps impeding the use of machine learning (ML)–based financial anomaly detection for village fund governance in Dana Desa. The following four analytical themes frame the literature: the Operationalization Chasm (no paper connects ML detection with the village governance context), the Scalability Illusion (detection AUC-ROC assumes conditions irrelevant to villages), the Absence of IS Theory (62\% of anomaly detection studies), and the Ground Truth Paradox (validation using artificial ground truth). A five-point dependency chain of structured gaps – three CRITICAL, one PARTIAL, and one METHODOLOGICAL – is identified and used to motivate a primary Design Science Research (DSR)-compliant study.

The central structural concept of the analysis is DSR × village = 0, across all 45 papers and for all sensitivity thresholds considered. While RELEVANCE (14 papers) and RIGOR (11 papers) knowledge exists in abundance, the artefact required to combine these two is missing. Closing of the three CRITICAL gaps (G1-G3) represents the design agenda of the proposed primary empirical work.

\textbf{Limitations}: include the lack of an SLR pre-registration protocol and lower DT5 scores for journals publishing studies related to Indonesia. For \textbf{Future research} directions, (1) DSR-compliant primary empirical work, using an unsupervised ensemble and IS Success Model, on anomaly detection using actual village funds disbursement data; (2) extending towards real-time anomaly monitoring using SiKPA/SIMDA village financial system data; (3) generalisation through repetition in other Indonesian provinces, and internationally in other decentralised finance contexts.

%% ─────────────────────────────────────────────────────────────────────────────
\section*{Acknowledgements}
This research was made possible by the support offered by Bina Nusantara University in the form of the Research and Publication Development Program. The authors express their sincere gratitude to the experience provided by experts from the Komisi Pemberantasan Korupsi (KPK), whose insight was instrumental in formulating the framework of governance adopted and corruption classification discussed in this paper. Finally, the authors would like to extend their thanks to the anonymous reviewers who provided valuable comments on earlier drafts. \textbf{AI Disclaimer:} AI has been used to improve the quality of writing.

%% =====================================================================
%% References — Vancouver numbered style, ordered by appearance
%% =====================================================================
\begin{thebibliography}{46}

\bibitem{method_kpk2023}
Komisi Pemberantasan Korupsi.
Laporan Monitoring Dana Desa 2015--2023.
KPK Directorate of Research and Development, 2023.
Jakarta, Indonesia.
URL: https://www.kpk.go.id/id/publikasi/laporan-kinerja

\bibitem{method_jensen1976}
Jensen, Michael C. and Meckling, William H.
Theory of the Firm: Managerial Behavior, Agency Costs and Ownership Structure.
\textit{Journal of Financial Economics}, 1976.
Volume 3, number 4, pages 305--360.
DOI: 10.1016/0304-405X(76)90026-X

\bibitem{p003}
Jack Nicholls and Aditya Kuppa and Nhien‐An Le‐Khac.
Financial Cybercrime: A Comprehensive Survey of Deep Learning Approaches to Tackle the Evolving Financial Crime Landscape.
\textit{IEEE Access}, 2021.
DOI: 10.1109/access.2021.3134076

\bibitem{p007}
Abdulwahab Ali Almazroi and Nasir Ayub.
Online Payment Fraud Detection Model Using Machine Learning Techniques.
\textit{IEEE Access}, 2023.
DOI: 10.1109/access.2023.3339226

\bibitem{p015}
José Edson de Albuquerque Filho and Laislla Carolina Pinheiro Brandão and Bruno Fernandes.
A Review of Neural Networks for Anomaly Detection.
\textit{IEEE Access}, 2022.
DOI: 10.1109/access.2022.3216007

\bibitem{p022}
Hanyao Gao and Gang Kou and Haiming Liang.
Machine learning in business and finance: a literature review and research opportunities.
\textit{Financial Innovation}, 2024.
DOI: 10.1186/s40854-024-00629-z

\bibitem{p026}
Shimal Sh. Taher and Siddeeq Y. Ameen and Jihan A. Ahmed.
Advanced Fraud Detection in Blockchain Transactions: An Ensemble Learning and Explainable AI Approach.
\textit{Engineering Technology \& Applied Science Research}, 2024.
DOI: 10.48084/etasr.6641

\bibitem{p041}
Elshan Gadimov and Ermiyas Birihanu.
Real-time suspicious detection framework for financial data streams.
\textit{International Journal of Information Technology}, 2025.
DOI: 10.1007/s41870-025-02529-6

\bibitem{p043}
Lucia Siciliani and Vincenzo Taccardi and Pierpaolo Basile.
AI-based decision support system for public procurement.
\textit{Information Systems}, 2023.
DOI: 10.1016/j.is.2023.102284

\bibitem{method_hevner2004}
Hevner, Alan R. and March, Salvatore T. and Park, Jinsoo and Ram, Sudha.
Design Science in Information Systems Research.
\textit{MIS Quarterly}, 2004.
Volume 28, number 1, pages 75--105.
DOI: 10.2307/25148625

\bibitem{p001}
Abdulalem Ali and Shukor Abd Razak and Siti Hajar Othman.
Financial Fraud Detection Based on Machine Learning: A Systematic Literature Review.
\textit{Applied Sciences}, 2022.
DOI: 10.3390/app12199637

\bibitem{p004}
Deepa Bisht and Rajesh Singh and Anita Gehlot.
Imperative Role of Integrating Digitalization in the Firms Finance: A Technological Perspective.
\textit{Electronics}, 2022.
DOI: 10.3390/electronics11193252

\bibitem{p037}
Luiz Moura and André Barcauí and Renan Carriço Payer.
AI and Financial Fraud Prevention: Mapping the Trends and Challenges Through a Bibliometric Lens.
\textit{Journal of risk and financial management}, 2025.
DOI: 10.3390/jrfm18060323

\bibitem{p008}
Francisco Cantú.
The Fingerprints of Fraud: Evidence from Mexico’s 1988 Presidential Election.
\textit{American Political Science Review}, 2019.
DOI: 10.1017/s0003055419000285

\bibitem{method_delone2003}
DeLone, William H. and McLean, Ephraim R.
The {DeLone} and {McLean} Model of Information Systems Success: A Ten-Year Update.
\textit{Journal of Management Information Systems}, 2003.
Volume 19, number 4, pages 9--30.
DOI: 10.1080/07421222.2003.11045748

\bibitem{method_prisma2020}
Page, Matthew J. and McKenzie, Joanne E. and Bossuyt, Patrick M. and Boutron, Isabelle and Hoffmann, Tammy C. and Mulrow, Cynthia D. and Shamseer, Larissa and Tetzlaff, Jennifer M. and Akl, Elie A. and Brennan, Sue E.
The {PRISMA} 2020 Statement: An Updated Guideline for Reporting Systematic Reviews.
\textit{BMJ}, 2021.
Volume 372, n71.
DOI: 10.1136/bmj.n71

\bibitem{method_thomas2008}
Thomas, James and Harden, Angela.
Methods for the Thematic Synthesis of Qualitative Research in Systematic Reviews.
\textit{BMC Medical Research Methodology}, 2008.
Volume 8, number 1, pages 45.
DOI: 10.1186/1471-2288-8-45

\bibitem{method_donthu2021}
Donthu, Naveen and Kumar, Satish and Mukherjee, Debmalya and Pandey, Nitesh and Lim, Weng Marc.
How to Conduct a Bibliometric Analysis: An Overview and Guidelines.
\textit{Journal of Business Research}, 2021.
Volume 133, pages 285--296.
DOI: 10.1016/j.jbusres.2021.04.070

\bibitem{method_templier2018}
Templier, Mathieu and {Par\'e}, Guy.
Transparency in Literature Reviews: An Assessment of Reporting Practices across Review Types and Genres in Top {IS} Journals.
\textit{European Journal of Information Systems}, 2018.
Volume 27, number 5, pages 503--550.
DOI: 10.1080/0960085X.2017.1398880

\bibitem{p002}
Josh Kamps and Bennett Kleinberg.
To the moon: defining and detecting cryptocurrency pump-and-dumps.
\textit{Crime Science}, 2018.
DOI: 10.1186/s40163-018-0093-5

\bibitem{p005}
Marc Schmitt.
Securing the digital world: Protecting smart infrastructures and digital industries with artificial intelligence (AI)-enabled malware and intrusion detection.
\textit{Journal of Industrial Information Integration}, 2023.
DOI: 10.1016/j.jii.2023.100520

\bibitem{p006}
Khyati Kapadiya and Usha Patel and Rajesh Gupta.
Blockchain and AI-Empowered Healthcare Insurance Fraud Detection: an Analysis, Architecture, and Future Prospects.
\textit{IEEE Access}, 2022.
DOI: 10.1109/access.2022.3194569

\bibitem{p009}
Samuel Oladiipo Olabanji and Yewande Alice Marquis and Chinasa Susan Adigwe.
AI-Driven Cloud Security: Examining the Impact of User Behavior Analysis on Threat Detection.
\textit{Asian Journal of Research in Computer Science}, 2024.
DOI: 10.9734/ajrcos/2024/v17i3424

\bibitem{p010}
Joanna Zhao and Xinruo Wang.
Unleashing efficiency and insights: Exploring the potential applications and challenges of ChatGPT in accounting.
\textit{Journal of Corporate Accounting \& Finance}, 2023.
DOI: 10.1002/jcaf.22663

\bibitem{p011}
Olubusola Odeyemi and Chinwe Chinazo Okoye and Onyeka Chrisanctus Ofodile.
INTEGRATING AI WITH BLOCKCHAIN FOR ENHANCED FINANCIAL SERVICES SECURITY.
\textit{Finance \& Accounting Research Journal}, 2024.
DOI: 10.51594/farj.v6i3.855

\bibitem{p012}
Odunayo Adewunmi Adelekan and Olawale Adisa and Bamidele Segun Ilugbusi.
EVOLVING TAX COMPLIANCE IN THE DIGITAL ERA: A COMPARATIVE ANALYSIS OF AI-DRIVEN MODELS AND BLOCKCHAIN TECHNOLOGY IN U.S. TAX ADMINISTRATION.
\textit{Computer Science \& IT Research Journal}, 2024.
DOI: 10.51594/csitrj.v5i2.759

\bibitem{p013}
Faozi A. Almaqtari.
The Role of IT Governance in the Integration of AI in Accounting and Auditing Operations.
\textit{Economies}, 2024.
DOI: 10.3390/economies12080199

\bibitem{p014}
Daniel Fährmann and Laura Martín and Luı́s Sánchez.
Anomaly Detection in Smart Environments: A Comprehensive Survey.
\textit{IEEE Access}, 2024.
DOI: 10.1109/access.2024.3395051

\bibitem{p016}
Felix Adebayo and Bakare Haslam and Olumide Ikumapayi.
Securing Government Revenue: A Cloud-Based AI Model for Predictive Detection of Tax-Related Financial Crimes.
\textit{International Journal of Computer Applications Technology and Research}, 2025.
DOI: 10.7753/ijcatr1405.1007

\bibitem{p017}
Davy Preuveneers and Vera Rimmer and Ilias Tsingenopoulos.
Chained Anomaly Detection Models for Federated Learning: An Intrusion Detection Case Study.
\textit{Applied Sciences}, 2018.
DOI: 10.3390/app8122663

\bibitem{p018}
K. Venkatesan and Syarifah Bahiyah Rahayu.
Blockchain security enhancement: an approach towards hybrid consensus algorithms and machine learning techniques.
\textit{Scientific Reports}, 2024.
DOI: 10.1038/s41598-024-51578-7

\bibitem{p019}
Mohammed Almehdhar and Abdullatif Albaseer and Muhammad Asif Khan.
Deep Learning in the Fast Lane: A Survey on Advanced Intrusion Detection Systems for Intelligent Vehicle Networks.
\textit{IEEE Open Journal of Vehicular Technology}, 2024.
DOI: 10.1109/ojvt.2024.3422253

\bibitem{p020}
Amal A. Alahmadi and Malak Aljabri and Fahd Alhaidari.
DDoS Attack Detection in IoT-Based Networks Using Machine Learning Models: A Survey and Research Directions.
\textit{Electronics}, 2023.
DOI: 10.3390/electronics12143103

\bibitem{p021}
George Kornaros.
Hardware-Assisted Machine Learning in Resource-Constrained IoT Environments for Security: Review and Future Prospective.
\textit{IEEE Access}, 2022.
DOI: 10.1109/access.2022.3179047

\bibitem{p023}
Darko Vuković and Senanu Dekpo-Adza and Stefana Matović.
AI integration in financial services: a systematic review of trends and regulatory challenges.
\textit{Humanities and Social Sciences Communications}, 2025.
DOI: 10.1057/s41599-025-04850-8

\bibitem{p024}
Yehia Ibrahim Alzoubi and Alok Mishra and Ahmet E. Topcu.
Research trends in deep learning and machine learning for cloud computing security.
\textit{Artificial Intelligence Review}, 2024.
DOI: 10.1007/s10462-024-10776-5

\bibitem{p025}
Ishaani Priyadarshini.
Anomaly Detection of IoT Cyberattacks in Smart Cities Using Federated Learning and Split Learning.
\textit{Big Data and Cognitive Computing}, 2024.
DOI: 10.3390/bdcc8030021

\bibitem{p027}
Saru Kumari.
Machine Learning Applications in Cryptocurrency: Detection, Prediction, and Behavioral Analysis of Bitcoin Market and Scam Activities in the USA.
\textit{International Journal of Sustainable Science and Technology}, 2025.
DOI: 10.22399/ijsusat.8

\bibitem{p028}
Sabina Szymoniak and Jacek Piątkowski and Mirosław Kurkowski.
Defense and Security Mechanisms in the Internet of Things: A Review.
\textit{Applied Sciences}, 2025.
DOI: 10.3390/app15020499

\bibitem{p029}
Tymoteusz Miller and Irmina Durlik and Ewelina Kostecka.
Artificial Intelligence in Maritime Cybersecurity: A Systematic Review of AI-Driven Threat Detection and Risk Mitigation Strategies.
\textit{Electronics}, 2025.
DOI: 10.3390/electronics14091844

\bibitem{p030}
Manuel J. C. S. Reis.
Edge-FLGuard: A Federated Learning Framework for Real-Time Anomaly Detection in 5G-Enabled IoT Ecosystems.
\textit{Applied Sciences}, 2025.
DOI: 10.3390/app15126452

\bibitem{p031}
Manuel J. C. S. Reis.
AI-Driven Anomaly Detection for Securing IoT Devices in 5G-Enabled Smart Cities.
\textit{Electronics}, 2025.
DOI: 10.3390/electronics14122492

\bibitem{p032}
Mueen Uddin and Muhammad Irshad and Irfan Ali Kandhro.
Generative AI revolution in cybersecurity: a comprehensive review of threat intelligence and operations.
\textit{Artificial Intelligence Review}, 2025.
DOI: 10.1007/s10462-025-11219-5

\bibitem{p033}
Olasehinde Omolayo and Ajao Ebenezer Taiwo and Tope David Aduloju.
Quantum Machine Learning Algorithms for Real-Time Epidemic Surveillance and Health Policy Simulation: A Review of Emerging Frameworks and Implementation Challenges.
\textit{International Journal of Multidisciplinary Research and Growth Evaluation}, 2024.
DOI: 10.54660/.ijmrge.2024.5.3.1100-1108

\bibitem{p034}
Mohammed M. Alwakeel.
AI-Assisted Real-Time Monitoring of Infectious Diseases in Urban Areas.
\textit{Mathematics}, 2025.
DOI: 10.3390/math13121911

\bibitem{p035}
Babawale Patrick Okare and Olasehinde Omolayo and Tope David Aduloju.
Designing Unified Compliance Intelligence Models for Scalable Risk Detection and Prevention in SME Financial Platforms.
\textit{International Journal of Multidisciplinary Research and Growth Evaluation}, 2025.
DOI: 10.54660/.ijmrge.2024.5.4.1421-1433

\bibitem{p036}
Arjun Kumar Bose Arnob and Rajarshi Roy Chowdhury and Nusrat Alam Chaiti.
A comprehensive systematic review of intrusion detection systems: emerging techniques, challenges, and future research directions.
\textit{Journal of Edge Computing}, 2025.
DOI: 10.55056/jec.885

\bibitem{p038}
Hamid Mirshekali and Mohammad Reza Shadi and Fatemehsadat Ghanadi Ladani.
A Review of Large Language Models for Energy Systems: Applications, Challenges, and Future Prospects.
\textit{IEEE Access}, 2025.
DOI: 10.1109/access.2025.3610994

\bibitem{p039}
Elizabeth Kuukua Amoako and Victor Boateng and Ola Ajay.
Exploring the role of Machine Learning and Deep Learning in Anti-Money Laundering (AML) strategies within U.S. Financial Industry: A systematic review of implementation, effectiveness, and challenges.
\textit{Finance \& Accounting Research Journal}, 2025.
DOI: 10.51594/farj.v7i1.1808

\bibitem{p040}
Andres J. Aparcana-Tasayco and Xianjun Deng and Jong Hyuk Park.
A systematic review of anomaly detection in IoT security: towards quantum machine learning approach.
\textit{EPJ Quantum Technology}, 2025.
DOI: 10.1140/epjqt/s40507-025-00414-6

\bibitem{p042}
Abhirup Khanna and Sapna Jain and Alessandro Burgio.
Blockchain-Enabled Supply Chain platform for Indian Dairy Industry: Safety and Traceability.
\textit{Foods}, 2022.
DOI: 10.3390/foods11172716

\bibitem{p044}
Beatrice Oyinkansola Adelakun and Ebere Ruth Onwubuariri and Gbenga Adeniyi Adeniran.
Enhancing fraud detection in accounting through AI: Techniques and case studies.
\textit{Finance \& Accounting Research Journal}, 2024.
DOI: 10.51594/farj.v6i6.1232

\bibitem{p045}
Xinze Yang and Chunkai Zhang and Yizhi Sun.
FinChain-BERT: A High-Accuracy Automatic Fraud Detection Model Based on NLP Methods for Financial Scenarios.
\textit{Information}, 2023.
DOI: 10.3390/info14090499

\bibitem{method_higgins2011}
Higgins, Julian P. T. and Green, Sally (editors).
Cochrane Handbook for Systematic Reviews of Interventions.
Edition 5.1.0.
The Cochrane Collaboration, 2011.
URL: https://handbook-5-1.cochrane.org/

\end{thebibliography}

\end{document}